% Previous Limitations section retained for reference.
% Our experimental results show that \sys{} generalizes across models and can accurately predict the energy consumption of parallelized inference on different hardware. One limitation in terms of the hardware is that we need to train \sys{} for each hardware as GPU capabilities and architecture can differ significantly. We note that, on the other hand, we do {\em not} have to train a new \sys{} predictor for each model. Our immediate next steps is to bridge this gap so that \sys{} becomes hardware-agnostic, further increasing its practical applicability across diverse deployment scenarios. Through our experiments, we have found hardware-agnostic energy prediction for LLMs to be a challenging task, warranting a full project in itself.
% \sys{} currently supports only decoder-style transformer models. %Extending it to encoder-based or encoder-decoder models is left for future work. 
% Encoder and encoder-decoder models often have bidirectional attention patterns and different execution flows, which affect both the structure of the model tree and the communication patterns. We believe that our parallelization-aware decomposition and communication modeling can be extended to encoder and encoder-decoder models; but we leave this extension for future work. Finally, in this work we show the effectiveness of \sys{} prediction under three popular parallelization strategies. Our next research steps will include extending the prediction to combinations of parallelism strategies such as 3D parallelism as well as other strategies including expert parallelism.
% 
% 
% \begin{comment}
% \begin{itemize}
%     \item only one type of model: vicuna
%     \item only focused on tensor parallelism, not pipeline parallel
%     \item parallel, not distributed
%     \item requires many experiments to collect significant sample size for wait times
%     \item can we call out some further opportunities to improve MAPE, especially for AllReduce?
% \end{itemize}
% \end{comment}

% \review{A next step is to reduce the cost of adapting \sys{} to new GPU platforms. The model-tree decomposition and structural features are reused across A6000, B40, and H200, while communication latency, power, and module-level energy predictors are calibrated for each platform. We will investigate whether limited target-platform profiling can support adaptation without retraining the predictors from scratch.}

\review{Our experimental results show that \sys{} generalizes across model families and accurately predicts the energy of parallelized inference on the evaluated GPU platforms.
%However, \sys{} is trained separately for each hardware as GPU capabilities and architecture can differ significantly.
While the model-tree decomposition and structural features are reused across A6000, B40, and H200, communication latency, power, and module-level energy predictors are calibrated separately for each platform.
A natural next step is to investigate whether limited target-platform profiling can support adaptation without retraining the predictors from scratch.}

% \review{We also plan to extend the model-tree abstraction to additional architectures and parallelism configurations. Encoder-decoder and mixture-of-experts models introduce different module structures and communication operations, while hybrid tensor, pipeline, and data parallelism combines multiple communication patterns. These settings call for corresponding extensions to module decomposition and communication-event calibration.}

\review{We also plan to extend the \sys{} methodology to additional architectures and parallelism configurations.
Encoder-decoder and mixture-of-experts models introduce different module structures and communication operations, while hybrid configurations combining tensor, pipeline, and data parallelism introduce multiple communication patterns.
Predicting energy in these settings will require extending the model-tree decomposition and characterizing the associated computation, communication, and synchronization costs.}

\review{\noindent\textbf{FSDP-style parameter sharding.} Fully Sharded Data Parallel (FSDP), originally developed for training, distributes parameter shards across ranks and all-gathers a module's weights before its forward computation~\cite{zhao2023fsdp}. Extending \sys{} to parameter-sharded inference requires modeling the volume and frequency of weight transfers under different sharding, prefetching, and parameter-retention policies. These transfers can overlap computation, and their energy cost must be distinguished from gradient synchronization, which is specific to training.}

\review{\noindent\textbf{Context parallelism.} Context-parallel methods distribute long sequences across GPUs, using schemes such as ring exchanges of attention key/value blocks~\cite{liu2023ringattention} or all-to-all redistribution of attention tensors~\cite{jacobs2023deepspeedulysses}. Predicting inference energy under these strategies requires modeling context-length-dependent attention work, per-GPU key/value storage, and communication schedules. Their potential overlap between attention computation and data transfers introduces additional energy-attribution challenges.}
